%!TEX root=../mythesis.tex
% Chapter 1

\chapter{Introduction} % Main chapter title
\chaptermark{Introduction}
\label{ch:introduction} % For referencing the chapter elsewhere, use \ref{Chapter1} 


%----------------------------------------------------------------------------------------
%	SECTION 1
%----------------------------------------------------------------------------------------

In the era of data-driven computing, geo-distributed service systems have emerged as a prevalent architecture for both large-scale technology enterprises and research institutions. To optimize data locality, these systems typically retain and process data at its generation site, which inherently requires dispatching the tasks of computation jobs across multiple geographically separated locations, depending on the spatial distribution of the requisite data chunks. By bridging compute and storage across wide-area networks (WANs), this architectural paradigm significantly enhances overall fault tolerance, reduces data access latency and minimizes cross-site bandwidth consumption \cite{hogade2022survey, bergui2021survey}. Consequently, this paradigm has expanded from big data analytics \cite{wu2017towards} and global cloud infrastructures (e.g., Google Cloud, Amazon Web Services) to a broader spectrum of applications, most notably machine learning \cite{hsieh2017gaia} and large-scale neural network model training \cite{tang2024fusionllm, wang2025spantrain}.

Despite these architectural advantages, operating job execution frameworks across geographically distributed sites inherently introduces a suite of new challenges, particularly for improving job processing efficiency and providing fair resource allocation. On one hand, from the user's perspective, tasks should be processed as expeditiously as possible to achieve the minimum job response time (the time span from the job arrival to the job completion). On the other hand, from the system's perspective, favoring certain jobs with extra resources naturally inflicts a performance penalty on others. This dichotomy creates an inherent conflict between efficiency and fairness. Furthermore, fairness in instantaneous resource allocation does not necessarily translate to long-term fairness in job response times, highlighting that the appropriate notion of fairness varies depending on the specific problem context. Additionally, finding the most efficient scheduling order in geo-distributed systems is known to be NP-hard, and the problem of ensuring fair allocation in such scenarios has not yet been widely studied. More importantly, the available bandwidth between different sites can be utilized to transfer tasks with required data chunks to improve either efficiency or fairness. However, how to properly schedule these network resources to achieve both objectives simultaneously introduces significant complexity to the problem. 

%We study the following three problems under the distributed job execution case:
This dissertation focuses on tackling the following three problems within the context of distributed job execution:
\begin{itemize}[nosep, leftmargin=*]
    \item \textbf{Problem 1: [Achieving fair instantaneous resource allocation with task transmission].} 
    How do we achieve fair, instantaneous resource allocation across jobs running concurrently in the system, given that the processing time of each task is unknown and only a limited number of tasks can be transferred based on the network bandwidth constraint in the system?
    
    \item \textbf{Problem 2: [Achieving efficient job execution with task transmission].} 
    How do we co-optimize the job scheduling and task transmission to maximize job processing efficiency, in scenarios where we are aware of the processing time of each task and a limited number of tasks can be transferred based on the network bandwidth constraint in the system?
    
    \item \textbf{Problem 3: [Achieving long-term fairness and efficiency trade-off].} 
    How do we schedule jobs so that the job processing efficiency is enhanced, while simultaneously providing a fair performance bound (i.e., completion deadlines) guaranteed for every job?
\end{itemize}

%We shall study and solve these problems in this thesis.
%This thesis studies and addresses the problems listed above.


\section{Research Scope}
\label{sec:scope}
In this thesis, we focus on the fairness and efficiency of distributed job execution with network resources and study three particular problems: \textit{max-min fair resource allocation with task transmission}, \textit{efficient job scheduling with task transmission}, and \textit{protective scheduling for balancing the long-term fairness and efficiency trade-off}.

The first part of the thesis is concerned with fairness in distributed job execution case that tasks (i.e., job resource demands) can be transferred across sites with limited network resources. Providing fair resource allocation can prevent starvation and fulfill Quality of Service (QoS) requirements for all jobs, but it is a non-trivial problem in distributed job execution. In addition, transferring tasks across the distributed system can potentially improve fairness since it enables starving jobs to access resources from other sites. However, effectively leveraging this requires sophisticated mechanism design to determine optimal job selection, source site origination, and target site placement decisions. In this part, we focus on achieving max-min fair single resource allocation with task transmission in distributed job execution, where the network resource is modeled by an upper bound on the transmission demand at each site. The problem consists of two critical hurdles: determining the distributed resource allocation for each job at each site with varying resource demand distributions, and utilizing the limited network resources to transfer the demands of certain jobs across sites to achieve more fair allocation and better overall utility. We aim to present a novel fairness notion and examine whether the notion satisfies widely-accepted fairness properties. We also aim to present algorithms of computing the resource allocation and task transmissions satisfying the fairness. %We present a resource allocation policy named Dynamic Aggregate Max-min Fairness (DAMF), which requires the aggregate resource allocation (i.e., the sum of resource allocation across all sites for each job) to be max-min fair among all possible allocations under all possible demand transmissions. Furthermore, we prove that DAMF satisfies widely-accepted fairness properties, including Pareto Efficiency, envy-freeness and strategy-proofness. We further present an algorithm to compute the exact DAMF resource allocation and a fast heuristic approximation to it. Experimental results show that DAMF outperforms the baseline policies in balancing resource allocation and enhancing job processing efficiency.

The second part of this thesis focuses on developing job scheduling algorithms which utilize network resources to improve job processing efficiency in geo-distributed job execution. In such scenarios, the completion time of each job is determined by the completion time of its last task across all sites. Therefore, the scheduling of jobs within individual sites presents significant optimization opportunities, especially when the task distributions of jobs among different sites are imbalanced. How to achieve optimal job scheduling in distributed job execution remains an NP-hard problem. Furthermore, transferring tasks (along with their requisite data) to other sites can potentially accelerate the execution of tasks and thereby improve the job's response time. However, there exists complex interdependency between the execution order of jobs and the transmission of tasks. Even the transmission of a single task may change the optimal execution order of jobs, while different job execution orders may, in turn, require different task transmissions to minimize the average job response time. We present a novel job scheduling strategy, called Workload-Aware Selection with Transmission (WAST), designed to optimize average job response time. WAST computes the global execution order of jobs and a set of task transmissions between different sites in an integrated manner, and iteratively schedules the job with the minimum response time until all jobs are scheduled. We also present a heuristic approximation algorithm that reduces the computational complexity without significantly compromising scheduling quality. Experimental results show that both algorithms can efficiently compute near-optimal job execution order and task transmissions, thereby significantly reducing job response times, especially in scenarios where the initial computation load distributions of jobs are highly skewed across sites. %Different from the existing work, we systematically analyze the synergistic impact of job scheduling and task transmission on job response time. We first focus on optimizing the response time of a single job, formulating the problem as a sequence of Integer Linear Programming (ILP) problems. Since solving ILPs is known to be NP-hard in the worst case, we further reduce the problem to a Minimum-cost Flow Problem that can be solved with polynomial-time complexity.

The third topic of this thesis is concerned with the long-term fairness and efficiency trade-off in the context of distributed job execution. In this topic, we aim to provide a performance guarantee—specifically, a completion time deadline for each job—and design scheduling policies that improve the job processing efficiency based on it. Inspired by the notion of protective scheduling in the single site case, our long-term fairness criterion dictates that no job should complete later than it would under a distributed Processor Sharing (PS) protocol, which evenly shares the site capacity among all active jobs at each site. Building on this foundation, we formulate the protective scheduling problem under the distributed job execution by establishing the slack system with two types of views, a global view and an independent view. Based on these two views, we present a set of algorithms extended from the single site scenario. We prove that although violations can be rare, no globally order-based scheduling policy is protective due to dynamic job arrivals and load distribution. To further enhance the efficiency, we propose a novel heuristic algorithm framework PFSP\_x, which not only guarantees protectiveness but also improves job processing efficiency compared to extension algorithms. We validate the proposed algorithms through real-world traces-driven evaluations. The results show that PFSP using the global slack scheme can effectively decrease the number of deadline violations while maintaining performance comparable to efficiency-targeted policies in terms of average job response time, and PFSP\_x achieves the highest job processing efficiency among all protective policies.
%we develop a job scheduling algorithm framework PFSP\_x that tackles the dual challenges of enhancing processing efficiency while strictly adhering to this fairness criterion. The main contributions of this work are summarized as follows:

\section{Thesis Contributions}
The remainder of this thesis is organized as follows:

In Chapter \ref{ch:literature_review}, we review the related work of fair resource allocation and efficient job scheduling in single-site and distributed job execution.

In Chapter \ref{ch:DAMF}, we study achieving max-min fair resource allocation with task (demand) transmission for distributed job execution. We present a fairness notion named Dynamic Aggregate Max-min Fairness (DAMF), which requires the aggregate resource allocation of each job to be max-min fair under all possible allocations and demand transmissions. We prove that DAMF satisfies important fairness properties, including Pareto Efficiency, envy-freeness and strategy-proofness. We further present an algorithm to search for DAMF allocation and a fast approximation method. The experiments using real-world job traces demonstrate that DAMF can effectively reduce the imbalance of the resource allocation among jobs while improving the average job response time.

In Chapter \ref{ch:WAST}, we study improving the job processing efficiency by job scheduling and task transmission in the distributed job execution case. We start from optimizing the response time of a single job and model it as solving multiple Integer Linear Programming (ILP) problems. We then transform the ILP into a minimum-cost flow problem which can be efficiently solved in polynomial time. Next, we present the policy Workload-Aware Selection with Transmission (WAST), which computes the global execution order of jobs and a set of task transmissions across sites in an integrated manner. Specifically, WAST iteratively schedules the job with the minimum response time, updates its workload and task transmission to the system, repeats this process until all jobs are scheduled. We also present a fast approximation to WAST with lower complexity. Our experiments using real-world job traces show that WAST and its approximation significantly outperform the baseline policies in terms of average job response time.

In Chapter \ref{ch:PFSP}, we study the long-term fairness and efficiency trade-off in distributed job execution. We dictate that no job should complete later than it would under a distributed Processor Sharing (PS) protocol as the criterion of long-term fairness, and develop job scheduling policies that improve efficiency with respect to this deadline. We construct two types of slack systems—a global view and an independent view—that help identify whether any job violates its PS deadline. We also extend algorithms presented in the single site case based on the two types of slacks. We prove that no globally order-based scheduling policy can maintain this guarantee because of unpredictable job arrival and load distribution. To solve this dilemma, we present an algorithm framework named PFSP\_x, which both guarantees protectiveness and improves job processing efficiency. Experimental results demonstrate that PFSP\_x achieves the lowest average job response time among all protective policies.

In Chapter \ref{ch:Con}, we conclude this thesis and outline possible future research directions.

